\chapter{Schluss}
Die durchgeführte Analyse und Evaluation der Klassifikationsmodelle zeigte, 
dass die Vorhersage von Sternebewertungen auf Basis von Nutzerrezensionen 
zahlreiche Herausforderungen mit sich brachte. Insbesondere die subjektive 
Natur der Bewertungen führte dazu, dass Texte nicht immer dem objektiv passenden 
Rating zugeordnet wurden. So traten beispielsweise Reviews mit einem überwiegend 
positiven Inhalt auf, denen lediglich vier Sterne statt fünf zugewiesen wurden, 
oder negative Texte, die mit drei statt zwei Sternen bewertet wurden. Dieses 
Verhalten spiegelte sich in den Modellergebnissen wider und erklärte teilweise 
die moderaten Accuracy-Werte der Multiclass-Klassifikation.\\
\\
\noindent Ein wesentlicher Faktor, der die Sentimentanalyse erschwerte, war die 
sprachliche Komplexität, die eine generelle Herausforderung in der natürlichen 
Sprachverarbeitung (NLP) darstellt. Typische Probleme treten bei Negationen, 
Ironie und Sarkasmus auf, da Modelle die sentimentale Wirkung solcher 
sprachlichen Mittel nur schwer korrekt erfassen können. Ebenso stellen 
Polysemie und mehrdeutige Wörter ein grundsätzlichen Problem in NLP dar: Ein 
Wort wie „mad“ kann je nach Kontext sowohl negative („angry“) als auch positive 
(„mad fun“) Bedeutungen annehmen.\\
\\
\noindent Darüber hinaus erschwert die Multipolarität von Texten die Klassifikation, da 
einzelne Texte oft gleichzeitig positive und negative Aspekte enthalten. Auch 
die Kontextabhängigkeit einzelner Wörter, sprachliche Nuancen, Redewendungen 
und grammatikalische Komplexität gelten allgemein als große Herausforderungen 
für NLP-Modelle. Diese Aspekte führen dazu, dass selbst leistungsfähige Modelle 
Schwierigkeiten haben, die subjektive Bewertung von Texten korrekt vorherzusagen 
[4].\\
\\
\noindent Zusammenfassend zeigte sich, dass die moderate Performance der Modelle, 
insbesondere der mittleren Klassen, durch eine Kombination aus subjektiver 
Bewertung, Mehrdeutigkeit, Kontextabhängigkeit und sprachlichen Besonderheiten 
bedingt war. Die logistische Regression erzielte dabei insgesamt leicht bessere 
Ergebnisse als Naive Bayes, insbesondere bei der Differenzierung der mittleren 
Klassen, konnte jedoch nicht alle oben genannten Herausforderungen vollständig 
kompensieren.\\
\\
\noindent Die Arbeit verdeutlichte, dass für eine robuste Sentimentanalyse weiterführende 
Ansätze, wie Transformer-Modelle, erforderlich 
sein könnten, um Negationen, Ironie, Polysemie und subjektive Bewertungen besser 
zu erfassen. Dennoch kann selbst bei den leistungsfähigsten Modellen die 
Vorhersage durch die Subjektivität der Bewertungen verzerrt sein. So kann 
ein Text, der nach seinem Inhalt eher einer 4-Sterne-Bewertung 
entspricht, von der bewertenden Person fälschlicherweise mit 3 oder 5 Sternen 
eingestuft worden sein, wodurch auch das beste Modell die exakte Sternezahl 
nicht zuverlässig vorhersagen kann.